% !TEX root = ../b-rg.tex
\chapter{Questions, Commands and Exclamations}\label{ch:sentence-types}

This chapter describes the forms of questions, answers, commands and exclamations, and a few fixed formulas used in conversation. The verbal forms involved, the imperative and the subjunctive, are described in Chapter~\ref{ch:tam}, and the position of pronouns after an imperative in Section~\ref{sec:imperative-pronouns}.

\section{Polar questions}\label{sec:polar-questions}
A yes--no question has the word order of a statement, with no inversion. It is marked by the question particle \bl{ou}, which follows the constituent being questioned. When the whole clause is questioned, as is usual, \bl{ou} therefore stands at the end.
    \glsen{Tu a fait pascr ag cat ou?}{ti a fet ˈpax.r̩ ɛj kat u}{{\Second\Sg} have {make.\Pst.\Ptcp} {eat-\Inf} {to.\Def} cat {\Q}}{Have you fed the cat?}
The particle is optional. A question without it, marked by intonation alone, is a little more colloquial: \bl{Scey un hom acas (ou)?} `Is there anyone at home?'.

Where the question concerns a verb with a complement clause, \bl{ou} stands straight after that verb and before the clause. This is the usual order with an indirect question.
    \glsen{Tu sa ou yon la hamt?}{ti sa u jɔn la hamt}{{\Second\Sg} know {\Q} where {\Third\Sg.\F} dwell}{Do you know where she lives?}
With a `whether' clause, \bl{ou} may instead come at the end of the whole sentence, and this is the written norm: \bl{Vos saveð ouð y micrman sey yert ou?} `Do you know if the bakery is open?' (where \bl{ouð} is `whether'; Section~\ref{sec:indirect-questions}). Placing \bl{ou} right after \bl{saveð} is possible in speech.

A negative question is a negative statement with \bl{ou}: \bl{Tu la boðau noc ou?} `Didn't you call her?'. On answering such questions, see Section~\ref{sec:answers}.

\subsection{Alternative questions}
A question offering a choice joins the alternatives with \bl{ben} `or' and has no \bl{ou}. Adding the particle turns it into a yes--no question about the whole, `do you want tea or coffee, or not?'.
    \glsen{Vos veleð tag ben caffay cavir?}{vɔz veˈlɛθ tɛj bɛn kaˈfe kaˈvɪr}{{\Second\Pl} {will-\Second\Pl.\Sbjv} tea or coffee {get-\Inf}}{Would you like tea or coffee?}

\subsection{Tag questions}
A question seeking confirmation adds the tag \bl{oc no's tu?} `isn't it so?' (literally `this isn't, you?') to a statement, whatever its verb. The tag agrees with the addressee: \bl{oc no soð vos?} in polite address. Colloquially it is shortened to \bl{no stu?}, \bl{no sovos?}, with the fused inverted copula (Section~\ref{sec:subject-pronouns}).
    \glsen{Tu vien souscen, oc no's tu?}{ti vjɛn suˈtsɛn ɔk nɔz ti}{{\Second\Sg} come tomorrow {\Prox.\Sg} {\Neg=be} {\Second\Sg}}{You're coming tomorrow, aren't you?}

\section{Content questions}\label{sec:content-questions}
The question words are \bl{cu} `who', \bl{cal} `what, which', \bl{yon} `where', \bl{cant} and \bl{candon} `when', \bl{com} `how', \bl{alcot} `how much, how many', \bl{hour alcot} `how long' and \bl{cascif} or \bl{carraçon} `why'. \bl{Cuscon} `whoever' is used as an incredulous `who' (Section~\ref{sec:interrogative-pronouns}).

The question word, or a phrase containing it, comes first, and the finite verb follows it immediately, so that the subject is inverted, whether it is a noun or a pronoun (Chapter~\ref{ch:simple-clause}). A preposition is always fronted with its question word, as in \bl{A cu} `to whom', \bl{Con cu} `with whom', \bl{Por cu} `for whom' and \bl{A cal} `about what'. It is never left at the end, as it is in relative clauses (Section~\ref{sec:oblique-relatives}).
    \glsen{A cu dau tu y moneð?}{a ki do ti i moˈnɛθ}{to who {give-\Pst} {\Second\Sg} {\Def} money}{Who did you give the money to?}
    \glsen{Con cu eu tu all'odysse?}{kɔn ki aw ti ˌa.lo.dɪˈse}{with who {go.\Pst} {\Second\Sg} {to.\Def=cinema}}{Who did you go to the cinema with?}
    \glsen{A cal es tu pesant?}{a kal ɛz ti peˈzant}{to what be {\Second\Sg} {think-\Prs.\Ptcp}}{What are you thinking about?}
Since the subject inverts while an object clitic stays before the verb, a question word as object and as subject is told apart by order: \bl{Cu galau tu all'odysse?} `Whom did you accompany to the cinema?', but \bl{Cu te galau all'odysse?} `Who accompanied you to the cinema?'. A verb that takes a direct object in Borlish needs no preposition, as with \bl{Cu es tu attenent?} `Who are you waiting for?'. `Whose' is \bl{a cu}, with the possessive preposition (\bl{A cu es l'oc çattrin?} `Whose umbrella is this?'; Section~\ref{sec:possessives}).

The question word may be reinforced with \bl{diabr} `devil', `on earth': \bl{Cu diabr te dis l'oç?} `Who on earth told you that?'.

\subsection{Several question words}
Only one question word is fronted. The others stay in place.
    \glsen{Cu accatau cal?}{ki ˌa.kaˈto kal}{who {buy-\Pst} what}{Who bought what?}

\subsection{Question words in place}
The question word also stays in its normal position in an echo question, which asks for something just said to be repeated or confirmed. It is then stressed, and the clause has the order of a statement.
    \glsen{Tu vis cu?}{ti vɪz ˈki}{{\Second\Sg} {see.\Pst} who}{You saw WHO?}
It likewise follows the verb with \bl{sceir} `there be' and \bl{cair} in its impersonal uses, which put their subject after the verb (Section~\ref{sec:impersonal}): \bl{Scey cu?} `Who's there?', \bl{Cay l'hour alcot?} `What time is it?'.
% CHECK: E212 (round 22) asks whether a question word can be extracted from an
% embedded clause ('Who do you think she saw?'), which relative clauses cannot
% do (Section~\ref{sec:relative-islands}).

\section{Indirect questions}\label{sec:indirect-questions}
An indirect content question begins with its question word. Unlike a direct question, it has the subject before the verb, both with nouns and with pronouns.
    \glsen{Jo say noc yon ty sour laur.}{ʒo se nɔk jɔn ti sur lor}{{\First\Sg} know {\Neg} where {\Second\Sg.\Poss} sister work}{I don't know where your sister works.}
Inversion is possible, though not required, when the verb is an auxiliary or the copula: \bl{Jo m'entrou yon son ðe eut} or \bl{yon ðe son eut} `I wonder where they've gone', \bl{yon es y quarf} `where the station is'. With a lexical verb it is unnatural, except perhaps with a very long subject.

An indirect yes--no question is introduced by \bl{ou} `whether', which has the form \bl{ouð} before a vowel. `Or whether' is \bl{ben ou}. After a verb of knowing that is negated or questioned, and after \bl{s'entrovar} `wonder' and \bl{doutar} `doubt', the clause is in the subjunctive (Section~\ref{sec:subjunctive}).
    \glsen{Jo m'entrou ouð Anna lo saif, ben ou Marc lo saif.}{ʒo mɛnˈtru uθ aˈna lo sef bɛn u mark lo sef}{{\First\Sg} {\First\Sg.\Obl=wonder} whether Anna {\Third\Sg.\N.\Acc} {know.\Sbjv} or whether Marc {\Third\Sg.\N.\Acc} {know.\Sbjv}}{I wonder whether Anna knows, or whether Marc does.}
The verb may be replaced by the pro-verb \bl{oc fair}, in the subjunctive (\bl{\ldots{} ben ou Marc oc faç}), but repeating the verb is preferred when the pro-verb is no shorter.

\section{Answers}\label{sec:answers}
`Yes' is \bl{ja}, usually followed by a full or short answer: \bl{Ja, jo vay ig par oy} `Yes, I'm going there now'. A short answer keeps the verb, with its complement replaced. A modal keeps \bl{oc} `this' as its object (\bl{Tu sa najar ou? --- Ja, j'oc sa} `Can you swim? --- Yes, I can'), a copula or auxiliary is repeated with its subject inverted after \bl{ja} or \bl{tan} `also' (\bl{Le's acas? --- Ja es le!} `Is he home? --- Yes, he is!'; \bl{Tan es la!} `So is she!'), and other verbs are replaced by the pro-verb \bl{oc fair} `do so'.

A short answer naming a person uses the disjunctive pronoun: \bl{Mey} `Me, it was me'. The pro-verb \bl{j'oc fey} `I did (so)' instead focuses on the fact that it was done.
    \glsen{Cu rompoy l'oc vas? Mey, j'oc fey.}{ki rɔmˈpɔj lɔk vaz mi ʒɔk fi}{who {break-\Pst} {\Def=\Prox.\Sg} vase {\First\Sg.\Dsj} {\First\Sg=\Prox.\Sg} {do.\Pst}}{Who broke this vase? I did.}
A negative question is contradicted not with \bl{ja} but with \bl{vis} `of course, indeed'. It may stand alone, be reinforced by \bl{certan} `certain', \bl{lon digr} `needless to say', \bl{clar} or, more formally, \bl{quippe} (\bl{Tu l'am noc ou? --- Vis, certan!} `Don't you like it? --- Of course I do!'), or be followed by the pro-verb, fronted with inversion.
    \glsen{Tu la boðau noc ou? Vis oc fey jo!}{ti la boˈðo nɔk u vɪz ɔk fi ʒo}{{\Second\Sg} {\Third\Sg.\F.\Acc} {call-\Pst} {\Neg} {\Q} indeed {\Prox.\Sg} {do.\Pst} {\First\Sg}}{Didn't you call her? Yes, I did!}
`No' as an answer is \bl{no}, the same word as the preverbal negator (Section~\ref{sec:neg-no}).
% CHECK: the word for 'maybe' as an answer has not been elicited (parcanç
% 'perhaps' is attested inside clauses).

\section{Commands and requests}\label{sec:commands-requests}
The forms of commands are described in Section~\ref{sec:commands}. In summary:
\begin{itemize}
    \item A command to one person uses the imperative, which is the present indicative singular (\bl{Vien ac!} `Come here!'). Five verbs use the subjunctive instead: \bl{sey}, \bl{ay}, \bl{faç}, \bl{saif}, \bl{remain}.
    \item A command to several people, or a polite one, uses the second person plural of the present subjunctive.
    \item `Let's' is the first person plural of the subjunctive (\bl{Fayau repaðr!} `Let's go home!').
    \item No subject pronoun is used, and object pronouns follow the verb (\bl{Da lo mey!} `Give it to me!').
\end{itemize}
    \glsen{Entreð eð asseyað, vos y dou!}{ɛnˈtrɛθ eθ ˌa.siˈjaθ vɔz i du}{{enter-\Second\Pl.\Sbjv} and {sit-\Second\Pl.\Sbjv} {\Second\Pl.\Dsj} {\Def} two}{Come in and sit down, both of you!}
A negative command places \bl{noc} after the verb, or uses \bl{skig} `avoid' with an infinitive (Section~\ref{sec:neg-noc}). A negative suggestion is the first person plural with \bl{noc}: \bl{Kerelleu noc dessur} `Let's not argue about it'.

Borlish has no word for `please'. A polite command uses \bl{veleð}, the second person plural subjunctive of \bl{volir}, with an infinitive, `be willing to', and the vocative follows it.
    \glsen{Veleð asseïr, domn.}{veˈlɛθ ˌa.siˈjɪr ˈdɔ.mn̩}{{will-\Second\Pl.\Sbjv} {sit-\Inf} lady}{Please sit down, madam.}
A request may also be made as a question, with the subjunctive of \bl{volir} and the final \bl{ou}: \bl{Vos veleð me digr ou yon es y quarf?} `Could you tell me where the station is?'. Between intimates, the plain indicative \bl{tu vel} is enough: \bl{Carin, tu vel y sal me passar?} `Darling, could you pass me the salt?'.

An offer, `let me', is \bl{ig} with the first person subjunctive, like French \textit{que je\ldots}, or the imperative of \bl{parmettr} `permit', or \bl{vel jo} with an infinitive (Section~\ref{sec:commands}).
    \glsen{Ig jo t'aïð!}{aj ʒo tɛˈjɪθ}{{\Comp} {\First\Sg} {\Second\Sg.\Obl=help.\Sbjv}}{Let me help you!}

\section{Exclamations}\label{sec:exclamatives}
Exclamations of degree are introduced by \bl{alcot} `how much'. With a noun it means `what (a)\ldots!', and the noun phrase follows, indefinite or bare, with any adjective in its usual place.
    \glsen{Alcot un bel majon!}{alˈkɔt ɪn bɛl maˈʒɔn}{how.much {\Indf} beautiful house}{What a beautiful house!}
Compare \bl{Alcot canottin douç!} `What a sweet little puppy!'. With a clause, \bl{alcot} causes inversion like a question word, and an adjective usually stays at the end. Fronting the adjective with \bl{alcot} (\bl{Alcot haut es tu denuð!}) is mannered.
    \glsen{Alcot es tu denuð haut!}{alˈkɔt ɛz ti deˈnɪθ hot}{how.much be {\Second\Sg} {become.\Pst.\Ptcp} tall}{How tall you've grown!}
In speech the copula and subject may be fused (\bl{Alcot stu denuð haut!}). A verb can be the focus as well: \bl{Alcot a tu te grandið!} `How you've grown!'.

\section{Formulas}\label{sec:formulas}
\begin{itemize}
    \item \textbf{Greetings.} The neutral greeting is \bl{bonjorn}. Slangier ones are \bl{hau}, \bl{hello} and \bl{yera}.
    \item \textbf{Address.} A vocative follows the clause, set off by a comma: \bl{domn} `madam', and politely \bl{bossir} `sir' and \bl{bedon} `madam' (from \bl{bon segr} and \bl{bel domn}), which are mostly used in address. Between intimates, \bl{carin} `darling'. A formal letter opens \bl{Bon Segr} `Dear Sir' (Chapter~\ref{ch:register-names}).
    \item \textbf{Thanks and courtesies.} `Thank you' is \bl{Benaul}, and `that's very kind' is \bl{de tal mensc} `of such kindness', one of the few expressions in which the archaic \bl{mensc} survives. `After you' is \bl{Leur vostr} `(the way is) free, yours'.
    \item \textbf{Other formulas.} \bl{Farau} `never mind, no worries'; \bl{Gevou!} `Careful!, Watch out!'; and the presentational \bl{Vey} `look, see', the imperative of \bl{veir}, used to point something out (\bl{Vey l'hom oç\ldots} `That's the man\ldots').
\end{itemize}
% CHECK: apologies and 'excuse me' (E214, round 22); toasts (siempr! in the
% plan); 'goodbye'.
